\begin{beginnerstatement}[Kurz erklärt: technische Abnahme]

Eine technische Abnahme prüft anhand festgelegter Schritte, ob eine
Softwarebasis die vorgesehenen technischen Anforderungen erfüllt. Ein lokaler
Nachweis entsteht auf einem Arbeitsrechner; ein Remote-Nachweis stammt etwa
von einem externen CI-System. Commitgebunden bedeutet, dass das Ergebnis nur
für den dabei geprüften Quellcodestand gilt. Die Abnahme für
\texttt{1cdfb6e} bestand mit offenen externen Punkten und bestätigt die
Eignung als Fallstudien-Baseline. Sie macht spätere Fehler auf
\texttt{a4487ac} nicht zu bestandenen Remote-Prüfungen. Produktionsreife würde
zusätzlich belegen, dass ein Produkt sicher und zuverlässig im echten Betrieb
eingesetzt werden kann. Diese Produktionsreife oder ein externes Deployment
wird durch die Abnahme ausdrücklich nicht bestätigt.

\end{beginnerstatement}
